\title{Controlling Text-to-Image Diffusion by Orthogonal Finetuning}

\begin{abstract}
State-of-the-art text-to-image diffusion models demonstrate remarkable proficiency in synthesizing photorealistic visuals from textual descriptions. A significant open question, however, is how to direct or modulate these advanced models to address a variety of downstream applications. In this work, we present a theoretically grounded approach termed Orthogonal Finetuning (OFT) for adapting text-to-image diffusion models to specific tasks. Distinct from prior techniques, OFT provably maintains hyperspherical energy, an attribute that captures pairwise neuron correlations when projected onto the unit hypersphere. Our analysis suggests that retaining this property is essential for upholding the semantic generative prowess of text-to-image diffusion models. To further ensure stable adaptation, we introduce Constrained Orthogonal Finetuning (COFT), which applies an extra radius constraint on the hypersphere. Our study focuses on two core finetuning tasks in the text-to-image domain: subject-driven generation, aiming to synthesize new images of a specific subject (given a handful of images and a prompt), and controllable generation, where models incorporate auxiliary control cues. Through empirical investigation, we demonstrate that our OFT framework yields superior results compared to existing alternatives, both in terms of output quality and convergence efficiency.
\end{abstract}

In this context, $\bm{W}$ stands for the OFT-adapted weight matrix, and $\bm{I}$ denotes the identity matrix. An overview of OFT is depicted in Figure~\ref{fig:block}. Much like the use of zero initialization in LoRA, it is necessary for OFT to adjust the pretrained model beginning from $\bm{W}^0$. To accomplish this requirement, we set the orthogonal matrix $\bm{R}$ to the identity at initialization, ensuring that the adapted model's weights match those of the pretrained model at the outset. To preserve the orthogonal property of $\bm{R}$, various differentiable orthogonalization methods from \cite{liu2021orthogonal,lezcano2019cheap} can be leveraged. We will elaborate on ensuring orthogonality in a manner that is computationally efficient.

\begin{itemize}[leftmargin=*,nosep]

\subsection{Re-scaled Orthogonal Finetuning}

Addressing the second question, a straightforward approach to adjust neuron directions involves applying a collective rotation or reflection across all neurons within the same layer, which corresponds to employing an orthogonal transformation. Although using alternative angular transformations—such as rotating neurons by distinct angles—offers increased flexibility, we observe that orthogonal transformation strikes an optimal balance between adaptability and regularization. In addition, \cite{liu2021orthogonal} demonstrates that orthogonal transformations possess sufficient representational power for neural network training. To further substantiate our claim, we carry out an experiment assessing the regularizing effect of orthogonality. This involves a controllable generation setup as in ControlNet~\cite{zhang2023adding}, with corresponding results shown in Figure~\ref{fig:ortho}. The results indicate that our conventional OFT yields stable training dynamics and precise control upon completion of training (epoch 20). In contrast, OFT without enforcing the orthogonality condition fails to produce realistic outputs and achieves no effective control. This experiment underscores the crucial role of the orthogonality constraint within OFT.

\begin{equation}
    \footnotesize
    \bm{R} = \text{diag}(\bm{R}_1, \bm{R}_2, \cdots, \bm{R}_r) = \begin{bmatrix}
    \bm{R}_1 \in \text{O}(\frac{d}{r}) & & \\
    & \ddots & \\
    & & \bm{R}_r \in \text{O}(\frac{d}{r})
    \end{bmatrix} \in \text{O}(d)
\end{equation}

\subsection{General Framework}

\textbf{Why not minimize hyperspherical energy}. A significant distinction from \cite{liu2021orthogonal} is that our objective is not to minimize hyperspherical energy. In classification tasks, maintaining non-redundant neurons is advantageous; forcing the neurons to be evenly distributed on the hypersphere (which minimum hyperspherical energy enforces) does not necessarily benefit finetuning, as this could undermine the information acquired during pretraining.

the series converges under the operator norm), which permits repositioning the constraint $\thickmuskip=2mu \medmuskip=2mu\|\bm{R}-\bm{I}\|\leq\epsilon$ within the Cayley transform. This leads to an analogous condition, $\thickmuskip=2mu \medmuskip=2mu \|\bm{Q}-\bm{0}\|\leq\epsilon'$, where $\bm{0}$ stands for the zero matrix and $\epsilon'$ is a distinct, tunable error threshold from $\epsilon$. Enforcing the new limitation on $\bm{Q}$ can be straightforwardly accomplished through projected gradient descent. To ensure the orthogonal matrix $\bm{R}$ is initialized as the identity, we initialize $\bm{Q}$ to be the zero matrix. The COFT approach thus embodies two explicit restrictions: a bounded change in hyperspherical energy and a tightly regulated departure from the pretrained model. While the latter is commonly present implicitly in existing finetuning methods, COFT formalizes it as an explicit constraint. Empirically, although OFT performs robustly, COFT demonstrates superior finetuning stability due to the explicit deviation regularization. Figure~\ref{fig:eps_coft} illustrates how varying $\epsilon$ modulates COFT's performance. Notably, larger values of $\epsilon$ grant greater adaptation freedom, making COFT resemble OFT, whereas decreasing $\epsilon$ draws COFT's outputs closer to those of the original pretrained diffusion model.

Modern text-to-image diffusion architectures~\cite{saharia2022photorealistic,ramesh2022hierarchical,rombach2022high} have established new benchmarks in producing high-fidelity images guided by text prompts. While these models exhibit robust performance, relying solely on textual instructions often leads to ambiguous outcomes, limiting precision and granular control in the generated imagery. In response, this work investigates two primary text-to-image generation scenarios:

\textbf{Balancing flexibility and regularization in finetuning.} Our observations reveal a fundamental trade-off between flexibility and regularity during finetuning. While standard finetuning methods allow for the highest level of flexibility, they are prone to instability and often result in model collapse. Remarkably, OFT employs a straightforward approach to achieve an effective equilibrium between flexibility and regularity by maintaining pairwise angles among neurons. The block-diagonal parameterization serves to further enhance the regularization effect on the orthogonal matrix.

\textbf{Convergence}. To assess convergence rates, we compare ControlNet, T2I-Adapter, LoRA, and COFT using the L2F task through both quantitative and qualitative analyses. Our quantitative approach involves calculating the average $\ell_2$ distance between the predicted facial landmarks and those from the control input. Figure~\ref{fig:face_convergence} presents the error curves for facial landmarks as a function of the number of finetuning epochs for each model. The plots clearly demonstrate that COFT and OFT attain substantially faster convergence relative to the baselines. Notably, LoRA requires 20 epochs to achieve the performance that OFT obtains by epoch 8. It is important to highlight that both OFT and COFT have a comparable count of trainable parameters to LoRA (and significantly fewer than ControlNet), yet their convergence rates considerably surpass those of existing approaches. The superior convergence efficiency of OFT is further evidenced in Figure~\ref{fig:teaser}: in the example on the right, OFT matches or exceeds the performance of ControlNet and LoRA while utilizing only 5\% of the ADE20K training set, attesting to its data efficiency. 

\subsection{Subject-driven Generation}

\textbf{OFT's architectural flexibility}. OFT is, in essence, applicable to any neural architecture. In the context of Transformers, LoRA is routinely used to adapt attention weights~\cite{hulora2022}, and thus, for a fair LoRA comparison, we also confine OFT to only updating attention weights in our experiments. Beyond dense layers, OFT is equally appropriate for convolutional layers; its block-diagonal $\bm{R}$ structure lends itself to interesting interpretations unique to convolution, not shared by LoRA. Specifically, when the number of blocks matches the input channel count, each block independently modifies a single neuron channel, analogous to depth-wise convolutions~\cite{chollet2017xception}. If all blocks in $\bm{R}$ are tied, a single orthogonal matrix is shared among all channels. An initial exploration of employing OFT to finetune convolutional layers can be found in Appendix~\ref{app:convolution}.

Here, $\thickmuskip=2mu \medmuskip=2mu \hat{\bm{w}}_i:=\bm{w}_i/\|\bm{w}_i\|$ designates the normalized $i$-th neuron, and the hyperspherical energy for $\bm{W}$ is specified by $\thickmuskip=2mu \medmuskip=2mu \text{HE}(\bm{W}):= \sum_{i\neq j} \|\hat{\bm{w}}_i-\hat{\bm{w}}_j\|^{-1}$. Notably, the optimal (minimum) value achievable in Eq.~\eqref{eq:oft_principle} is zero, which occurs if and only if the relation between $\bm{W}$ and $\bm{W}^0$ is purely a rotation or a reflection, i.e., $\thickmuskip=2mu \medmuskip=2mu \bm{W}=\bm{R}\bm{W}^0$ for some orthogonal matrix $\thickmuskip=2mu \medmuskip=2mu \bm{R}\in\mathbb{R}^{d\times d}$ (where a determinant of $1$ indicates a rotation, and $-1$ indicates a reflection). This formulation is central to OFT, which restricts finetuning to learning orthogonal matrices shared across each layer to linearly

\subsection{Efficient Orthogonal Parameterization}\label{sect:eff_ortho}

\begin{equation}\label{eq:coft_general}
    \footnotesize
    \bm{z}=\bm{W}^\top\bm{x}=(\bm{R}\cdot\bm{W}^0)^\top\bm{x},~~~\text{s.t.}~\bm{R}^\top\bm{R}=\bm{R}\bm{R}^\top=\bm{I},~\norm{\bm{R}-\bm{I}}\leq \epsilon
\end{equation}

\section{Introduction}

Nonetheless, OFT brings forth several open questions. To enforce orthogonality, OFT utilizes Cayley parametrization, necessitating a matrix inversion that can impede scalability. Although adopting block diagonal parametrization partially mitigates this issue, efficiently performing differentiable matrix inversion remains a significant challenge. Furthermore, OFT exhibits promising compositional properties: the orthogonal matrices obtained from separate OFT finetuning processes can be composed through multiplication, with the product remaining orthogonal. Investigating whether these compositions effectively retain the learned information from all tasks is a compelling future research direction. Lastly, the parameter efficiency of OFT hinges on the block diagonal structure, which, while beneficial, may introduce systematic biases and constrain adaptability. Advancing the parameter efficiency with less bias and greater effectiveness stands out as an important unresolved question.

To address the first question, we are motivated by empirical findings in \cite{liu2018decoupled,chen2020angular}, which demonstrate that the angular difference between features effectively captures the semantic gap. SphereNet~\cite{liu2017deep} provides evidence that when all neurons in a neural network are normalized onto a unit hypersphere during training, the model maintains comparable capacity and often exhibits superior generalization, suggesting that the orientation of the neurons encapsulates the most critical information from the data. To further elucidate the significance of neuron angles, we present a toy example in Figure~\ref{fig:toy_ae}. Here, a conventional convolutional autoencoder is trained on a dataset of flower images. During training, the feature map is generated using the standard inner product (where $z$ represents the convolution kernel output from $\bm{w}$ over an input $\bm{x}$ within a sliding window). For the test phase, we evaluate three strategies for feature map construction: (a) employing the same inner product as during training, (b) utilizing only the feature magnitudes, and (c) considering only angular (directional) information. The outcomes, displayed in Figure~\ref{fig:toy_ae}, reveal that the angular components of the neurons can almost flawlessly reconstruct the original images, whereas the magnitude conveys little to no informative content. It is worth noting that no cosine-based activation (as in (c)) was involved in the training process—the model was trained exclusively with the inner product. These observations indicate that neuron directions predominantly encode the semantic details of the input images. Consequently, fine-tuning neuron angles is likely to be a particularly effective strategy for modifying semantic content in images.

\textbf{Finetuning stability and convergence}. We begin by assessing the robustness and convergence rate of DreamBooth, LoRA, OFT, and COFT during finetuning. As illustrated in Figure~\ref{fig:teaser} and Figure~\ref{fig:db_convergence}, both COFT and OFT demonstrate strong stability in adapting the diffusion model. Within 400 iterations, DreamBooth and OFT-based approaches successfully capture the desired content, whereas LoRA fails to maintain subject identity. By the 2000th iteration, DreamBooth starts producing degraded, collapsed outputs, and LoRA generates yellow fur in place of a yellow shirt. Conversely, OFT and COFT continue to exhibit reliable and coherent image control at this stage. These findings affirm that the OFT framework combines rapid convergence with stable subject-driven generation. Furthermore, this stability with respect to iteration count is highly advantageous, as it reduces the burden of extensive hyperparameter tuning across different subjects. For OFT and COFT, one can safely select a large iteration count without meticulous adjustment. In addition, when using COFT with an appropriately chosen $\epsilon$, both the learning rate and iteration count become straightforward to set without the need for careful optimization.

\textbf{Controllable generation}. Image-to-image translation falls within the broader category of controllable generation tasks. Historically, conditional generative adversarial networks have been the predominant approach to this problem~\cite{isola2017image,zhu2017toward,wang2018high,park2019semantic,choi2018stargan}. More recently, diffusion models have been leveraged for image-to-image translation as well~\cite{saharia2022palette,voynov2022sketch,wang2022pretraining}. Innovations such as ControlNet~\cite{zhang2023adding} have demonstrated effective strategies for controlling pretrained diffusion models through finetuning and integrating additional control inputs, thus achieving impressive performance in controllable generation. A related line of work, T2I-Adapter~\cite{mou2023t2i}, also explores the finetuning of pretrained diffusion models to improve control over the output images. In alignment with these research efforts~\cite{zhang2023adding,mou2023t2i}, we employ OFT to finetune pretrained diffusion models and consistently attain enhanced controllability with both reduced training data and fewer finetuning parameters. Importantly, OFT incurs no extra computational burden during inference.

\textbf{Why OFT converges faster}. Figure~\ref{fig:toy_ae} demonstrates that modifying neuron orientation is most effective for altering semantic content—a change that OFT is particularly suited to enact. Furthermore, OFT can be interpreted as performing finetuning on a smooth hyperspherical manifold, which leads to a more favorable optimization landscape. This perspective is further substantiated by empirical results in \cite{liu2021orthogonal}.

\subsection{Constrained Orthogonal Finetuning}

\textbf{Settings}. We adopt ControlNet~\cite{zhang2023adding}, T2I-Adapter~\cite{mou2023t2i}, and LoRA~\cite{hulora2022} as comparison baselines. The experiments encompass three complex controllable generation benchmarks: Canny edge to image~(C2I) using the COCO dataset~\cite{lin2014microsoft}, segmentation map to image~(S2I) on ADE20K~\cite{zhou2017scene}, and landmark to face~(L2F) with CelebA-HQ~\cite{karras2018progressive,xia2021tedigan}. All approaches perform finetuning on Stable Diffusion~(SD) v1.5 for 20 epochs on each dataset. Extended results can be found in Appendix~\ref{app:more_qualitative},~\ref{app:more_control_task}, and~\ref{app:failure}.

\begin{itemize}[leftmargin=*,nosep]

\textbf{Subject-driven generation}. Addressing the challenge of subject modification, prior methods~\cite{avrahami2022blended,nichol2022glide} introduce user-provided masks as supplementary inputs. Inversion-based techniques~\cite{choi2021ilvr,dhariwal2021diffusion,ramesh2022hierarchical,gal2022image} offer means to alter the scene context while maintaining subject identity. \cite{hertz2022prompt} achieves both local and global edits without requiring masks. Nonetheless, these strategies often struggle to retain detailed identity features of the subject. Pivotal Tuning~\cite{roich2022pivotal} instead refines a generator in the vicinity of an inverted latent representation, aided by a regularization term to better preserve identity. Similarly,~\cite{nitzan2022mystyle} learns a custom facial prior from several portraits of an individual. \cite{casanova2021instance} produces diverse versions of an instance, albeit sometimes at the expense of instance-specific details. DreamBooth~\cite{ruiz2023dreambooth} adapts text-to-image diffusion models by leveraging a personalized token and a few subject images, optimizing with a reconstruction loss plus a class-specific prior preservation objective. While adopting the general setup of DreamBooth, OFT distinguishes itself by employing orthogonal transformations in the finetuning process, rather than relying solely on standard finetuning with reduced learning rates.

\section{Experiments and Results}

\textbf{General experimental setup.} For all experiments, Stable Diffusion v1.5~\cite{rombach2022high} serves as the pretrained backbone for text-to-image generation. To maintain fairness, samples generated by different methods are randomly selected. We follow the standard DreamBooth~\cite{ruiz2023dreambooth} procedure for subject-driven generation tasks, while the protocols from ControlNet~\cite{zhang2023adding} and T2I-Adapter~\cite{mou2023t2i} are adopted for controllable generation tasks. For an unbiased comparison with LoRA, OFT or COFT is applied only at the layers where LoRA operates. Additional experimental results and implementation specifics can be found in Appendix~\ref{app:settings}.

\textbf{Quantitative comparison}. To benchmark performance quantitatively, and consistent with the evaluation in \cite{ruiz2023dreambooth}, we measure subject fidelity using DINO~\cite{caron2021emerging} and CLIP-I~\cite{radford2021learning}, text prompt alignment with CLIP-T~\cite{radford2021learning}, and generation diversity using LPIPS~\cite{zhang2018unreasonable}. Specifically, CLIP-I evaluates the mean cosine similarity between CLIP embeddings from real and generated images, while DINO provides an analogous measure but uses ViT S/16 DINO features. CLIP-T reports the average cosine similarity between generated image and text prompt embeddings in CLIP space. We also analyze LPIPS by calculating the average pairwise cosine similarity for outputs generated from the same prompt and subject. As reported in Table~\ref{table:dreambooth_final}, both COFT and OFT surpass DreamBooth and LoRA on the DINO and CLIP-I metrics by a significant margin, and perform similarly or better on prompt fidelity and diversity. Each method was tested by finetuning the base model 30 times with different random seeds to ensure consistency. Overall, these experiments demonstrate that OFT not only enables faster and more stable convergence, but also delivers consistently superior final outcomes.

Furthermore, our approach draws on the concept of orthogonal over-parameterized training~\cite{liu2021orthogonal}, which entails training classification neural networks by applying orthogonal transformations to a randomly initialized architecture. In such initializations, hyperspherical energy is expected to be low, and the objective in \cite{liu2021orthogonal} is to sustain this low energy throughout training, as reduced energy correlates with enhanced generalization~\cite{liu2018learning,lin2020regularizing}. That work demonstrates that orthogonal transformations afford sufficient expressiveness for training well-generalizing classification models. By contrast, our work centers on the finetuning of text-to-image diffusion models to achieve improved controllability and superior generative outcomes on downstream tasks. We clarify how OFT differs from the approach in \cite{liu2021orthogonal} along two main lines. First, whereas \cite{liu2021orthogonal} is designed to lower hyperspherical energy, OFT is explicitly crafted to maintain the same hyperspherical energy as that of the pretrained model, thereby preserving the inherent structure and semantics established during pretraining—a consideration especially vital for diffusion models where minimizing energy can be detrimental. Second, OFT incorporates an explicit regularization to minimize departure from the pretrained parameters, motivating a constrained formulation; in contrast, \cite{liu2021orthogonal} operates without such constraints. Achieving a balance between adaptability and preservation is essential in finetuning, and we contend that the OFT framework provides a robust mechanism for this purpose. Our main contributions are as follows:

\textbf{Experimental setup.} DreamBooth~\cite{ruiz2023dreambooth} and LoRA~\cite{hulora2022} serve as comparative baselines. All methods use the same loss formulation as outlined in DreamBooth. The training procedures and hyperparameters for DreamBooth and LoRA generally adhere to recommendations in their respective original papers, employing optimal configurations. Expanded results and further analysis appear in Appendices~\ref{app:settings},~\ref{app:coft_oft},~\ref{app:more_qualitative}, and~\ref{app:failure}.

To motivate the use of orthogonal transformation in finetuning text-to-image diffusion models, we break down the discussion into two key points: (1) the rationale for adjusting neuron angles (\emph{i.e.}, their directions) during finetuning, and (2) the justification for utilizing orthogonal transformation specifically for tuning these angles.

For qualitative analysis, we concentrate on comparing OFT, COFT, and ControlNet, given that ControlNet’s landmark error is most closely aligned with OFT. Visualizations in Figure~\ref{fig:face_convergence_qual} illustrate that both OFT and COFT exhibit stable convergence, with synthesized face poses progressively aligning with the intended control landmarks over training. Conversely, ControlNet displays abrupt improvements in controllability after epoch 8, a sudden convergence phenomenon also reported by \cite{zhang2023adding}. The smooth and stable convergence trajectory of OFT notably simplifies practical deployment, as its training behavior is more predictable and reliable, thereby facilitating the search for suitable hyperparameters.

    \item \textbf{Subject-driven generation}~\cite{ruiz2023dreambooth}: Starting from a small set of images depicting a specific subject, this task involves producing novel renditions of the same subject embedded within varied contexts, utilizing additional text cues.
    \item \textbf{Controllable generation}~\cite{zhang2023adding,mou2023t2i}: This setting introduces auxiliary control information (\emph{e.g.}, canny edge maps, segmentation layouts), challenging the model to generate outputs that adhere to these controls in conjunction with a textual prompt.
\end{itemize}

For illustration, consider a fully connected layer described as $\thickmuskip=2mu \medmuskip=2mu \bm{W}=\{\bm{w}_1,\cdots,\bm{w}_n\}\in\mathbb{R}^{d\times n}$, with each column vector $\bm{w}_i\in\mathbb{R}^{d}$ representing the $i$-th neuron and $\bm{W}^0$ denoting its pretrained counterpart. The output $\thickmuskip=2mu \medmuskip=2mu \bm{z}\in\mathbb{R}^n$ is produced by $\thickmuskip=2mu \medmuskip=2mu \bm{z}=\bm{W}^\top\bm{x}$ for input $\thickmuskip=2mu \medmuskip=2mu \bm{x}\in\mathbb{R}^d$. In this context, OFT may be regarded as encouraging the hyperspherical energy gap between the finetuned and original layers to be minimized:

Although classical methods such as the Gram-Schmidt process can maintain differentiability, they are computationally demanding and impractical for large-scale use~\cite{liu2021orthogonal}. For improved computational efficiency, we utilize the Cayley parameterization, which produces an orthogonal matrix as follows: $\thickmuskip=2mu \medmuskip=2mu \bm{R} = (\bm{I} + \bm{Q})(I - \bm{Q})^{-1}$, where $\bm{Q}$ is a skew-symmetric matrix, i.e., $\thickmuskip=2mu \medmuskip=2mu \bm{Q} = -\bm{Q}^\top$. This approach, while more efficient, imposes a mild constraint: the Cayley parameterization yields only special orthogonal matrices, i.e., with determinant 1. Empirically, we observe this constraint does not impact overall performance. Nevertheless, directly applying Cayley transforms to parameterize $\bm{R}$ remains parameter-intensive, especially for large values of $d$. To mitigate this, we decompose $\bm{R}$ into a block-diagonal matrix containing $r$ blocks, resulting in the structure:

    \item We introduce a new finetuning technique, Orthogonal Finetuning, designed to enhance controllability in text-to-image diffusion models. To boost stability, we also present a constrained version of this method, which restricts the model’s angular change from its pretrained state.
    \item In contrast to prior finetuning strategies, OFT executes model adaptation while rigorously maintaining hyperspherical energy. Our empirical analysis suggests that this energy is a key indicator for conserving the pretrained model’s generative semantics.
    \item We implement OFT for both subject-driven and controllable generation tasks, conducting an extensive empirical analysis that shows clear improvements in generation fidelity, convergence rate, and finetuning stability over earlier techniques. Additionally, OFT demonstrates greater sample efficiency, attaining effective convergence with far fewer training images and epochs.
    \item In the context of controllable generation, we propose a novel control consistency metric to measure controllability. This metric evaluates the consistency between the control signal inferred from generated samples and the original input signal, thereby providing further evidence for the robust controllability achieved by OFT.
\end{itemize}

\subsection{Controllable Generation}

\textbf{Qualitative comparison}. To provide a more immediate and visual sense of OFT’s advantages, we display several randomly selected cases for subject-driven generation in Figure~\ref{fig:dreambooth_final}. For equitable comparison, each example is derived from the identical finetuned model corresponding to each approach, and no individual text prompts have been specifically tuned to optimize final outcomes. The chosen model for each method corresponds to the highest validation CLIP scores obtained. As shown in Figure~\ref{fig:dreambooth_final}, both OFT and COFT maintain strong semantic consistency in subject representation, whereas LoRA often exhibits significant failure in preserving subject identity (\emph{e.g.}, in the bowl example, LoRA entirely fails to retain the subject’s identity). Additionally, OFT and COFT enable notably finer control via textual input; in contrast, DreamBooth, while capable of retaining subject identity, frequently does not produce images that align with the intended prompt (\emph{e.g.}, see the first row for the bowl example). The qualitative assessment indicates that OFT simultaneously excels in both control and subject preservation. Furthermore, the number of OFT iterations does not notably impact performance, allowing it to maintain strong results even with increased iteration counts—a property not shared by either DreamBooth or LoRA. Supplementary qualitative examples appear in Appendix~\ref{app:more_qualitative}. Additional validation of OFT’s superiority is provided through a human study in Appendix~\ref{app:human}.

\textbf{Connection to LoRA}. LoRA mitigates drastic alterations to the pretrained weight matrix by introducing an additional low-rank matrix. In contrast, OFT preserves pretraining knowledge by enforcing the transformation applied to the pretrained weights to be orthogonal (thus full-rank), thereby safeguarding against potential corruption of preexisting information. The forward computation in OFT can be expressed as $\thickmuskip=2mu \medmuskip=2mu \bm{z}=(\bm{R}\bm{W}^0)^\top\bm{x}=(\bm{W}^0+(\bm{R}-\bm{I})\bm{W}^0)^\top\bm{x}$, where $\thickmuskip=2mu \medmuskip=2mu (\bm{R}-\bm{I})\bm{W}^0$ functions similarly to LoRA's low-rank perturbation. Given that $\bm{W}^0$ generally possesses full rank, OFT also effects a low-rank adjustment if $\thickmuskip=2mu \medmuskip=2mu \bm{R}-\bm{I}$ is low-rank. Comparable to LoRA’s utilization of a rank parameter $r'$, OFT introduces a diagonal block parameter $r$ to constrain the number of learnable parameters. Intriguingly, LoRA and OFT reflect two fundamentally different strategies for parameter efficiency: LoRA leverages low-rank characteristics, while OFT capitalizes on structural sparsity via block-diagonal orthogonality.

The rationale behind employing orthogonal transformations in neuron finetuning can be partly traced to the 2D Fourier transform, which decomposes an image into magnitude and phase components. Of these, the phase spectrum—reflecting angular relationships between inputs and bases—encodes the principal semantic content. Notably, one can pair the phase spectrum of an image with a random magnitude spectrum to reconstruct an image that retains its semantics. This observation indicates that adjusting the orientation of neurons is fundamental to semantic alteration of generated images, as is required in both subject-driven and controllable image generation. However, allowing unrestricted adjustment of neuron orientations risks compromising the generative capability established during pretraining. To mitigate this, we advocate for maintaining the angles between all pairs of neurons, based on the supposition that such inter-neuron angles are integral to encoding the knowledge in neural networks. Following this reasoning, we introduce a layer-wise, shared orthogonal transformation for neurons, ensuring that hyperspherical energy is conserved across layers.

\subsection{Why Does Orthogonal Transformation Make Sense?}

At their core, both tasks address the challenge of how to adapt text-to-image diffusion models through finetuning while maintaining the generative capabilities established during pretraining. An effective finetuning approach should meet two primary criteria: (1) \emph{training efficiency}, meaning reduced numbers of trainable parameters and minimal training epochs, and (2) \emph{generalizability preservation}, which involves retaining the model's ability to generate high-quality and varied outputs. To achieve these aims, finetuning is commonly implemented by either incrementally adjusting existing neuron weights with a low learning rate (\emph{e.g.}, \cite{ruiz2023dreambooth}), or by augmenting the model with compact modules containing re-parameterized weights (\emph{e.g.}, \cite{hulora2022,zhang2023adding}). Although straightforward, both strategies are limited in their ability to consistently uphold the generative strengths imparted by pretraining. Furthermore, there is a scarcity of rigorous frameworks to guide the design of robust finetuning procedures or to select key hyperparameters, such as the optimal number of epochs. A significant challenge lies in the absence of quantitative measures that assess how much of the original generative capacity is retained. Many current finetuning approaches operate under the presumption that maintaining a smaller Euclidean distance between the updated and pretrained model weights correlates with better retention of the pretrained functions, which encourages the use of very conservative learning rates. However, while this principle sometimes aligns with empirical outcomes, we contend that simply minimizing the Euclidean distance does not adequately reflect semantic preservation. Therefore, adopting more structurally informed metrics for comparing finetuned and pretrained models has the potential to substantially improve the maintenance of pretraining capabilities and to enhance finetuning robustness.

\textbf{Quantitative comparison}. To assess controllable generation, we propose a control consistency metric, which measures how closely the control signal derived from the output image matches the original input control signal. For the C2I scenario, IoU and F1 score are calculated; in the S2I context, we report mean IoU along with mean and overall accuracy. For L2F, the evaluation is based on the average $\ell_2$ distance between the intended control landmarks and those predicted by the model. Detailed descriptions of these consistency metrics can be found in Appendix~\ref{app:settings}. All methods are evaluated using their optimal hyperparameter configurations. As reported in Table~\ref{table:control_gen}, OFT and COFT consistently demonstrate superior control and precision over alternative approaches. Adapter-based methods (such as T2I-Adapter and ControlNet) exhibit slower convergence rates and inferior ultimate performance. While LoRA surpasses ControlNet on the S2I task, its results fall behind on C2I and L2F tasks. Moreover, we observe that LoRA’s maximum achievable performance remains limited, regardless of intensive hyperparameter tuning. By contrast, the OFT framework’s performance continues to improve as the model size increases, without evident saturation. Importantly, our approach to quantitative evaluation for controllable generation serves as a notable contribution: it enables precise assessment of the control fidelity in fine-tuned models, within the accuracy bounds imposed by the pre-trained segmentation or detection models.

\begin{equation}\label{eq:oft_general}
    \footnotesize
    \bm{z} = \bm{W}^\top\bm{x} = (\bm{R} \cdot \bm{W}^0)^\top\bm{x},~~~\text{s.t.}~\bm{R}^\top\bm{R} = \bm{R}\bm{R}^\top = \bm{I}
\end{equation}

Drawing inspiration from the pivotal role that angular relationships between neurons play in shaping visual semantics, we introduce a straightforward yet robust approach—orthogonal finetuning—for exerting control over text-to-image diffusion models. This method is particularly tailored for two primary text-to-image tasks: subject-driven generation and controllable generation. Relative to existing approaches, OFT provides enhanced controllability and increased finetuning stability while utilizing a smaller set of trainable parameters. Additionally, OFT does not impose extra computational burden during inference, which supports efficient deployment in real-world scenarios.

Here, $\text{O}(d)$ refers to the group of all orthogonal matrices in $d$ dimensions, where $\thickmuskip=2mu \medmuskip=2mu \bm{R}\in\mathbb{R}^{d\times d}$ and each $\thickmuskip=2mu \medmuskip=2mu \bm{R}_i\in\mathbb{R}^{d/r\times d/r},\forall i$ represents an orthogonal matrix. In the special case when $\thickmuskip=2mu \medmuskip=2mu r=1$, the block-diagonal structure reduces to a full unconstrained orthogonal matrix. For a square orthogonal matrix of dimension $\thickmuskip=2mu \medmuskip=2mu d\times d$, the number of parameters required is $\thickmuskip=2mu \medmuskip=2mu d(d-1)/2$, corresponding to a computational complexity of $\mathcal{O}(d^2)$. In contrast, an orthogonal matrix partitioned into $r$ blocks has $\thickmuskip=2mu \medmuskip=2mu d(d/r-1)/2$ parameters, and therefore a reduced complexity of $\thickmuskip=2mu \medmuskip=2mu \mathcal{O}(d^2/r)$. Further efficiency can be achieved by parameter sharing, specifically by setting $\thickmuskip=2mu \medmuskip=2mu \bm{R}_i = \bm{R}_j$ for all $i \neq j$, which brings the parameter complexity down to $\mathcal{O}(d^2/r^2)$. Despite these parameter-reduction techniques, the overall matrix $\bm{R}$ remains orthogonal and thus maintains the ability to preserve hyperspherical energy.

Drawing on empirical findings that hyperspherical similarity effectively captures semantic relationships~\cite{liu2017deep,liu2018decoupled,chen2020angular}, we utilize the concept of hyperspherical energy~\cite{liu2018learning} to describe the pairwise relationships among neurons. Hyperspherical energy, computed as the aggregate hyperspherical similarity (\emph{e.g.}, cosine similarity) across all neuron pairs within a single layer, serves as a quantitative measure of how uniformly neurons are distributed over the unit hypersphere~\cite{liu2021learning}. Our working assumption is that optimal finetuning should induce minimal change in hyperspherical energy relative to the original pretrained model. A straightforward strategy might involve penalizing deviations in hyperspherical energy during finetuning; however, such an approach cannot guarantee a precise minimization of the difference. Instead, we exploit an important property of hyperspherical energy: it remains invariant under any shared orthogonal transformation applied to all neurons, thereby preserving pairwise hyperspherical similarities. Building on this invariance, we introduce Orthogonal Finetuning~(OFT), a method that adapts large text-to-image diffusion models to specific tasks while ensuring their hyperspherical energy is unaltered. This approach involves learning an orthogonal transformation common to each layer’s neurons so that their mutual angular relationships are maintained. Conceptually, OFT can also be regarded as realigning the canonical coordinate basis for neurons within the layer. Further, noting that closer Euclidean proximity between the finetuned and pretrained models is typically associated with greater retention of pretrained capabilities, we present an extension named Constrained Orthogonal Finetuning~(COFT), which restricts the adaptation to lie within a fixed-radius hypersphere centered at the pretrained neuron configuration.

\section{Concluding Remarks and Open Problems}

It is possible to further restrict the expressiveness of the OFT approach by ensuring the finetuned model stays within a bounded proximity to the pretrained solution. This leads to the COFT strategy, wherein the forward computation is described as:

\textbf{Text-to-image diffusion models}. The field of text-to-image generation has advanced rapidly~\cite{saharia2022photorealistic,ramesh2022hierarchical,rombach2022high,gu2022vector,nichol2022glide}, driven by improvements in diffusion-based generative modeling~\cite{ho2020denoising,song2021score,song2021denoising,dhariwal2021diffusion} as well as progress in multimodal vision-language representation learning~\cite{su2020vlbert,lu2019vilbert,radford2021learning,wang2021simvlm,li2022blip,li2023blip,alayrac2022flamingo,schuhmann2022laion}. In approaches such as GLIDE~\cite{nichol2022glide} and Imagen~\cite{saharia2022photorealistic}, diffusion models operate in pixel space. GLIDE employs paired text-image data to jointly train its text encoder alongside a diffusion prior, whereas Imagen uses a fixed, pretrained text encoder. In latent space approaches like Stable Diffusion~\cite{rombach2022high} and DALL-E2~\cite{ramesh2022hierarchical}, Stable Diffusion leverages VQ-GAN~\cite{esser2021taming} to create a visual codebook for latent representations, while DALL-E2 utilizes CLIP~\cite{radford2021learning} to build a shared latent embedding space uniting images and textual descriptions. Beyond diffusion architectures, text-to-image research has also explored generative adversarial networks~\cite{reed2016generative,zhang2017stackgan,xu2018attngan,li2019controllable} and autoregressive models~\cite{ramesh2021zero,ding2021cogview,wu2022nuwa,yuscaling}. OFT stands out as a model-agnostic approach to finetuning and is compatible with any text-to-image diffusion architecture.

\section{Intriguing Insights and Discussions}

\section{Orthogonal Finetuning}

\begin{equation}\label{eq:oft_principle}
\footnotesize
\min_{\bm{W}} \left\| \text{HE}(\bm{W})-\text{HE}(\bm{W}^0)\right\|~~\Leftrightarrow~~\min_{\bm{W}} \bigg{\|} \sum_{i\neq j} \|\hat{\bm{w}}_i-\hat{\bm{w}}_j\|^{-1}-\sum_{i\neq j} \|\hat{\bm{w}}^0_i-\hat{\bm{w}}^0_j\|^{-1}\bigg{\|}
\end{equation}

\textbf{Qualitative comparison}. We further provide a qualitative assessment of OFT and COFT in relation to established leading methods such as ControlNet, T2I-Adapter, and LoRA. As illustrated in Figure~\ref{fig:control_final}, randomly sampled outputs indicate that OFT and COFT consistently deliver images of superior realism and fidelity, in addition to offering precise control over generation. For the S2I task, LoRA fails to respect the input segmentation maps, whereas ControlNet, OFT, and COFT are all capable of accurately guiding image synthesis. Compared to ControlNet, OFT and COFT not only yield images with richer detail and more plausible geometrical structures, but do so with significantly fewer parameters. On the C2I task, OFT and COFT can generate semantically meaningful images even when provided with coarse Canny edge inputs, outperforming T2I-Adapter and LoRA by a wide margin. Regarding the L2F task, our approach achieves superior pose accuracy, producing highly controllable face generation even for difficult facial orientations. Across all three control scenarios, the qualitative outcomes from OFT and COFT surpass those from competing state-of-the-art baselines, highlighting the efficacy of our OFT framework in controllable image generation. For greater thoroughness, additional qualitative comparisons across all three control tasks are presented in Appendix~\ref{app:control_exp}, and we further demonstrate OFT’s generalization ability on diverse control tasks—including dense pose to human body, sketch to image, and depth to image translation—in Appendix~\ref{app:more_control_task}.

\textbf{Inference efficiency unchanged.} In contrast to ControlNet, our OFT approach does not increase inference-time computational demands for the finetuned model. At inference, the orthogonal matrix $\bm{R}$ learned during training can be directly incorporated with the original pretrained weight matrix $\bm{W}^0$ by simple matrix multiplication, giving the resultant matrix $\thickmuskip=2mu \medmuskip=2mu \bm{W}=\bm{R}\bm{W}^0$. This operation preserves inference speed, matching that of the base pretrained model.

\section{Related Work}

A comparison of OFT and LoRA in terms of parameter efficiency is now presented. For LoRA with a low-rank parameter $r'$, the total number of trainable parameters is $\thickmuskip=2mu \medmuskip=2mu r'(d+n)$. Assuming both $r$ and $r'$ scale linearly with the neuron dimension $d$ (for instance, letting $\thickmuskip=2mu \medmuskip=2mu r=r'=\alpha d$, where $\thickmuskip=2mu \medmuskip=2mu 0<\alpha\leq 1$ is some fixed constant), LoRA exhibits a parameter complexity of $\mathcal{O}(d^2+dn)$, whereas OFT can achieve $\mathcal{O}(d)$. To make this comparison more tangible, consider a weight matrix sized $\thickmuskip=2mu \medmuskip=2mu 128\times 128$; in this case, LoRA, using $\thickmuskip=2mu \medmuskip=2mu r'=8$, requires $2,048$ trainable parameters, while OFT, with $\thickmuskip=2mu \medmuskip=2mu r=8$ (and without employing block sharing), uses $960$ trainable parameters.

Finetuning in the conventional sense generally leverages gradient descent with a modest learning rate to modify either the whole model or selected layers. The reduced learning rate effectively penalizes significant deviations from the initial, pretrained model, thus conventional finetuning implicitly operates to minimize $\thickmuskip=2mu \medmuskip=2mu \| \bm{M} - \bm{M}^0\|$, where $\bm{M}$ corresponds to the updated model parameters and $\bm{M}^0$ refers to the original, pretrained parameters. While this implicit restriction is intuitively reasonable, it may lack sufficient rigidity when finetuning large-scale models. LoRA addresses this potential shortcoming by adding an explicit low-rank restriction on the parameter update: $\thickmuskip=2mu \medmuskip=2mu \text{rank}(\bm{M} - \bm{M}^0)=r'$, where $r'$ is a specified small integer. In contrast, OFT employs a different kind of constraint, focusing on preserving pairwise neuron similarities by enforcing $\thickmuskip=2mu \medmuskip=2mu \| \text{HE}(\bm{M})-\text{HE}(\bm{M}^0) \|=0$, where $\text{HE}(\cdot)$ captures the hyperspherical energy for the weight matrix. 

This framework imposes both an orthogonality requirement and an additional restriction, which limits the norm of the deviation from the identity matrix to $\epsilon$. While orthogonality can be enforced via the Cayley parameterization (refer to Section~\ref{sect:eff_ortho}), embedding the $\epsilon$-deviation constraint within this parameterization poses challenges. To better understand the behavior of the Cayley transform, we utilize the Neumann series to approximate $\thickmuskip=2mu \medmuskip=2mu \bm{R}=(\bm{I}+\bm

We introduce a straightforward augmentation to OFT by allowing a per-neuron scaling factor to be learned. The rationale is that rescaling the output of neurons does not alter hyperspherical energy, since this normalization negates their magnitude. In practice, we redefine the forward operation as $\thickmuskip=2mu \medmuskip=2mu\bm{z}=(\bm{R}\bm{W}^0\bm{D})^\top\bm{x}$\footnote{\textbf{Errata}: The NeurIPS camera ready version incorrectly wrote the re-scaled OFT forward path as $\thickmuskip=2mu \medmuskip=2mu\bm{z}=(\bm{D}\bm{R}\bm{W}^0)^\top\bm{x}$. The original codebase is correct, so the outcomes shown in Appendix~\ref{app:rescaled} remain valid.}, with $\thickmuskip=2mu \medmuskip=2mu \bm{D}=\text{diag}(s_1,\cdots,s_n)\in\mathbb{R}^{n\times n}$ representing a learnable diagonal matrix where each diagonal entry $s_1,\ldots,s_n$ is positive. Whereas standard OFT only updates $\bm{R}$ (see Eq.~\eqref{eq:oft_general}), this revised formulation allows both $\bm{R}$ and $\bm{D}$ to adapt during training. Re-scaled OFT thereby enhances OFT's expressive capacity at the cost of minimal parameter increase. Our main experiments focus on the vanilla OFT, in order to underscore the efficacy of orthogonal transforms alone, but we observe that the re-scaled OFT generally outperforms the basic variant (refer to Appendix~\ref{app:rescaled}).

\textbf{Model finetuning}. Adapting large pretrained models to new tasks via finetuning has become a mainstream strategy~\cite{devlin2018bert,brown2020language,he2022masked}. Within the context of finetuning, parameter-efficient adaptation techniques (\emph{e.g.}, \cite{hulora2022,houlsby2019parameter,pfeiffer2020adapterfusion}) are widely examined in the natural language processing domain. Among these, LoRA~\cite{hulora2022} shares the closest resemblance to OFT, as it is based on the premise of low-rank additive updates to weights during finetuning. OFT, by contrast, applies a multiplicative, layer-shared orthogonal transformation to modify neuron weights, which theoretically preserves pairwise angular relationships between neurons in the same layer, thereby ensuring improved stability.